% ======================================================================
% Chủ đề 2 -- Hình học và Đo lường -- phần c (bài 31--40)
% ======================================================================
\clearpage
\dividerbar{Mở rộng -- Nâng cao}

\begin{probbox}{31}{\relax}
Trong mặt phẳng toạ độ $Oxy$, cho điểm $A(-3;\,2)$.
\begin{enumerate}[label=\alph*),leftmargin=1.6em,itemsep=0pt,topsep=2pt]
  \item Gọi $B$ là điểm đối xứng với $A$ qua trục $Ox$, $C$ là điểm đối xứng với $A$ qua trục $Oy$, $D$ là điểm đối xứng với $A$ qua gốc $O$. Tìm toạ độ của $B$, $C$, $D$ và cho biết mỗi điểm nằm ở góc phần tư nào.
  \item Bốn điểm $A$, $C$, $D$, $B$ là bốn đỉnh của một hình chữ nhật. Tính chu vi và diện tích hình chữ nhật đó.
\end{enumerate}
\end{probbox}
\probsol{
\textbf{Phân tích đề.} ``Đối xứng qua một trục'' giống như soi gương qua trục đó: điểm ảnh nằm ở phía bên kia trục, cách trục một khoảng bằng điểm ban đầu. Ta xem toạ độ nào đổi dấu, toạ độ nào giữ nguyên.
\begin{recallbox}Đối xứng qua trục $Ox$: $(x;\,y)\to(x;\,-y)$. Đối xứng qua trục $Oy$: $(x;\,y)\to(-x;\,y)$. Đối xứng qua gốc $O$ (cũng là quay $180^\circ$ quanh $O$): $(x;\,y)\to(-x;\,-y)$.\end{recallbox}
\step{1}{Điểm $B$ -- đối xứng qua $Ox$}
Soi gương qua trục hoành: điểm giữ nguyên vị trí trái -- phải (hoành độ không đổi), chỉ đổi từ phía trên xuống phía dưới (tung độ đổi dấu). $B(-3;\,-2)$, dạng $(-;-)$: góc phần tư III.

\step{2}{Điểm $C$ -- đối xứng qua $Oy$}
Soi gương qua trục tung: tung độ giữ nguyên, hoành độ đổi dấu. $C(3;\,2)$, dạng $(+;+)$: góc phần tư I.

\step{3}{Điểm $D$ -- đối xứng qua gốc $O$}
Cả hai toạ độ đổi dấu: $D(3;\,-2)$, dạng $(+;-)$: góc phần tư IV.
\begin{gmFig}[scale=0.55]
  \fill[tmBlue!12] (-3,-2) rectangle (3,2);
  \draw[Line,line width=0.3pt] (-4.5,-3.5) grid (4.5,3.5);
  \draw[->,Navy,thick] (-4.8,0)--(4.9,0) node[right,text=Ink]{$x$};
  \draw[->,Navy,thick] (0,-3.8)--(0,3.9) node[above,text=Ink]{$y$};
  \foreach \t in {-4,...,-1,1,2,3,4}{\node[below left,font=\scriptsize,inner sep=1.5pt] at (\t,0){$\t$};}
  \foreach \t in {-3,...,-1,1,2,3}{\node[below left,font=\scriptsize,inner sep=1.5pt] at (0,\t){$\t$};}
  \node[below left,font=\scriptsize,inner sep=1.5pt] at (0,0){$O$};
  \draw[tmBlue,thick] (-3,-2) rectangle (3,2);
  \fill[AnsRed] (-3,2) circle (3pt) node[above left]{$A$};
  \fill[tmPurple] (-3,-2) circle (3pt) node[below left]{$B$};
  \fill[tmGreen!70!black] (3,2) circle (3pt) node[above right]{$C$};
  \fill[tmOrange] (3,-2) circle (3pt) node[below right]{$D$};
  \node[above,font=\footnotesize,text=tmBlue] at (1.5,2){$6$};
  \node[right,font=\footnotesize,text=tmBlue] at (3,1){$4$};
\end{gmFig}

\step{4}{Tính kích thước hình chữ nhật}
$A(-3;\,2)$ và $C(3;\,2)$ cùng tung độ: $AC=3-(-3)=6$.
$A(-3;\,2)$ và $B(-3;\,-2)$ cùng hoành độ: $AB=2-(-2)=4$.

\step{5}{Tính chu vi và diện tích}
\begin{recallbox}Hình chữ nhật dài $a$, rộng $b$: $P=2\cdot(a+b)$, $S=a\cdot b$.\end{recallbox}
$P=2\cdot(6+4)=20$ (đơn vị độ dài); $\;S=6\cdot4=24$ (đơn vị diện tích).
\begin{tipbox}Hình chữ nhật này nhận cả hai trục toạ độ làm trục đối xứng. Mỗi góc phần tư chứa đúng một phần tư hình, là hình chữ nhật $3\times2$ có diện tích $6$; và $4\cdot6=24$.\end{tipbox}
\verify{Mỗi góc phần tư chứa một điểm: $A$ (II), $C$ (I), $B$ (III), $D$ (IV), đúng như hình.}
\ans{a) $B(-3;\,-2)$: III; $C(3;\,2)$: I; $D(3;\,-2)$: IV. \quad b) Chu vi $20$, diện tích $24$.}
}

\begin{probbox}{32}{\relax}
Trong hình vẽ, tam giác $ABC$ có $\widehat{BAC}=40^\circ$; cạnh $BC$ được kéo dài về cả hai phía, cạnh $AB$ được kéo dài qua $B$. Góc ngoài tại đỉnh $C$ bằng $115^\circ$. Giá trị của $x$ là
\begin{gmFig}[scale=0.55]
  \coordinate (B) at (0,0); \coordinate (C) at (4,0); \coordinate (A) at (1.460,5.448);
  \coordinate (L) at (-1.6,0); \coordinate (D) at (-0.42,-1.567); \coordinate (E) at (5.8,0);
  \draw[Ink,thick] (L)--(E); \draw[Ink,thick] (D)--(A)--(C);
  \pic[draw,thick,angle radius=0.7cm,"$40^\circ$",angle eccentricity=1.45]{angle=B--A--C};
  \pic[draw,thick,angle radius=0.45cm,"$115^\circ$",angle eccentricity=2.1]{angle=E--C--A};
  \pic[draw,AnsRed,thick,angle radius=0.45cm,"$x^\circ$",angle eccentricity=1.8]{angle=L--B--D};
  \node[above] at (A){$A$}; \node[above left] at (B){$B$}; \node[below] at (C){$C$};
\end{gmFig}
\choices{$65$}{$75$}{$105$}{$115$}{$140$}
\end{probbox}
\probsol{
\textbf{Phân tích đề.} Góc $x^\circ$ ở đỉnh $B$ được tạo bởi hai tia kéo dài (tia đối của $BC$ và tia đối của $BA$), nên nó \emph{đối đỉnh} với góc trong $\widehat{ABC}$. Vậy chỉ cần tìm $\widehat{ABC}$. Ta giải bằng hai cách.

\step{1}{Cách 1 -- qua góc trong tại $C$}
\begin{recallbox}Hai góc kề bù có tổng $180^\circ$; tổng ba góc trong tam giác bằng $180^\circ$.\end{recallbox}
Góc trong tại $C$ kề bù với góc ngoài $115^\circ$: $\widehat{ACB}=180^\circ-115^\circ=65^\circ$.
Góc trong tại $B$: $\widehat{ABC}=180^\circ-40^\circ-65^\circ=75^\circ$.

\step{2}{Cách 2 -- dùng tính chất góc ngoài}
\begin{recallbox}Góc ngoài của tam giác bằng tổng hai góc trong không kề với nó.\end{recallbox}
Góc ngoài tại $C$ bằng $\widehat{A}+\widehat{B}$: $115^\circ=40^\circ+\widehat{ABC}$, suy ra $\widehat{ABC}=115^\circ-40^\circ=75^\circ$.

\step{3}{Suy ra $x$}
\begin{recallbox}Hai góc đối đỉnh thì bằng nhau.\end{recallbox}
$x^\circ=\widehat{ABC}=75^\circ$, nên $x=75$.
\begin{warnbox}
\begin{itemize}[leftmargin=1.3em,itemsep=0pt,topsep=0pt]
  \item (A) $65$ là góc trong tại $C$, không phải tại $B$.
  \item (C) $105$ là góc kề bù với $\widehat{ABC}$ -- nhầm ``kề bù'' với ``đối đỉnh''.
\end{itemize}
\end{warnbox}
\verify{$40^\circ+75^\circ+65^\circ=180^\circ$ và $40^\circ+75^\circ=115^\circ$ (góc ngoài tại $C$). Cả hai điều kiện đều đúng.}
Vậy chọn \textbf{(B)}.
\ans{(B) $75$}
}

\begin{probbox}{33}{\relax}
\begin{enumerate}[label=\alph*),leftmargin=1.6em,itemsep=0pt,topsep=0pt]
  \item Một tam giác cân có góc ở đỉnh gấp $4$ lần mỗi góc ở đáy. Tính số đo các góc của tam giác.
  \item Một tam giác cân có một góc bằng $36^\circ$. Tìm tất cả các khả năng cho số đo hai góc còn lại. Với mỗi khả năng, tính hiệu giữa góc lớn nhất và góc nhỏ nhất của tam giác.
  \item Có tam giác cân nào có một góc ở đáy bằng $92^\circ$ không? Vì sao?
\end{enumerate}
\end{probbox}
\probsol{
\textbf{Phân tích đề.} Cả ba câu đều dựa trên hai sự thật: tam giác cân có hai góc ở đáy bằng nhau, và tổng ba góc bằng $180^\circ$.
\begin{recallbox}Tam giác cân: góc ở đỉnh $d$, mỗi góc ở đáy $g$, thì $d+2g=180^\circ$.\end{recallbox}
\step{1}{Câu a -- gọi ẩn và lập phương trình}
Gọi số đo mỗi góc ở đáy là $x$ (độ), với $0<x<90$. Góc ở đỉnh gấp $4$ lần mỗi góc ở đáy nên có số đo $4x$.
Tổng ba góc của tam giác bằng $180^\circ$, nên ta có phương trình
\[ x+x+4x=180. \]
Thu gọn vế trái: $6x=180$. Chia cả hai vế cho $6$: $x=180:6=30$.
Vậy mỗi góc ở đáy bằng $30^\circ$, góc ở đỉnh bằng $4\cdot30^\circ=120^\circ$.

\step{2}{Câu b -- xét hai trường hợp}
\emph{Trường hợp 1: $36^\circ$ là góc ở đỉnh.} Mỗi góc đáy: $(180^\circ-36^\circ):2=144^\circ:2=72^\circ$. Tam giác $36^\circ,\,72^\circ,\,72^\circ$. Hiệu góc lớn nhất và nhỏ nhất: $72^\circ-36^\circ=36^\circ$.

\emph{Trường hợp 2: $36^\circ$ là góc ở đáy.} Góc đáy kia $36^\circ$, góc ở đỉnh $180^\circ-72^\circ=108^\circ$. Tam giác $36^\circ,\,36^\circ,\,108^\circ$. Hiệu: $108^\circ-36^\circ=72^\circ$.

\step{3}{Câu c -- lập luận}
\begin{proofbox}Nếu một góc ở đáy bằng $92^\circ$ thì góc ở đáy còn lại cũng bằng $92^\circ$. Khi đó riêng hai góc đáy đã có tổng $184^\circ$, lớn hơn $180^\circ$ -- vô lí, vì tổng \emph{ba} góc chỉ bằng $180^\circ$. Vậy không có tam giác cân nào như thế. (Tổng quát: góc ở đáy của tam giác cân luôn nhỏ hơn $90^\circ$.)\end{proofbox}
\verify{$30+30+120=180$; $\;36+72+72=180$; $\;36+36+108=180$.}
\ans{a) $30^\circ$, $30^\circ$, $120^\circ$. \ b) $72^\circ$ và $72^\circ$ (hiệu $36^\circ$), hoặc $36^\circ$ và $108^\circ$ (hiệu $72^\circ$). \ c) Không có.}
}

\begin{probbox}{34}{\stars{1}}
Trong mặt phẳng toạ độ, cho điểm $A(4;\,1)$.
\begin{enumerate}[label=\alph*),leftmargin=1.6em,itemsep=0pt,topsep=2pt]
  \item Quay $A$ quanh gốc $O$ một góc $90^\circ$ ngược chiều kim đồng hồ được điểm $B$; quay $B$ thêm $90^\circ$ ngược chiều kim đồng hồ được $C$; quay $C$ thêm $90^\circ$ ngược chiều kim đồng hồ được $E$. Tìm toạ độ $B$, $C$, $E$.
  \item Quay $A$ quanh $O$ một góc $270^\circ$ \emph{theo chiều kim đồng hồ} thì được điểm nào trong ba điểm trên?
  \item Tứ giác $ABCE$ là một hình vuông. Tính diện tích của nó.
\end{enumerate}
\end{probbox}
\probsol{
\textbf{Phân tích đề.} Ta dùng quy tắc quay $90^\circ$ trong mặt phẳng toạ độ, sau đó tính diện tích hình vuông ``nghiêng'' bằng cách đặt nó vào trong một hình vuông lớn có cạnh nằm theo các đường kẻ ô.
\begin{recallbox}Quay $90^\circ$ ngược chiều kim đồng hồ quanh $O$: $(x;\,y)\to(-y;\,x)$. Quay $\alpha$ theo chiều kim đồng hồ giống quay $360^\circ-\alpha$ ngược chiều.\end{recallbox}
\begin{proofbox}Vì sao $(x;\,y)\to(-y;\,x)$? Điểm $A(4;\,1)$ đạt được bằng cách đi ``sang phải $4$, lên $1$''. Quay cả đường đi này $90^\circ$ ngược chiều kim đồng hồ: ``sang phải'' biến thành ``lên trên'', ``lên trên'' biến thành ``sang trái''. Vậy ảnh đạt được bằng ``lên $4$, sang trái $1$'', tức điểm $(-1;\,4)$.\end{proofbox}
\step{1}{Câu a -- áp dụng quy tắc ba lần}
$A(4;\,1)\to B(-1;\,4)\to C(-4;\,-1)\to E(1;\,-4)$.

(Kiểm tra: $C$ là ảnh của $A$ qua phép quay $180^\circ$, đúng là $(-4;\,-1)$, cả hai toạ độ đổi dấu.)

\step{2}{Câu b}
Quay $270^\circ$ theo chiều kim đồng hồ giống quay $360^\circ-270^\circ=90^\circ$ ngược chiều kim đồng hồ, nên được đúng điểm $B(-1;\,4)$.

\needspace{14\baselineskip}
\step{3}{Câu c -- đặt hình vuông vào khung}
\begin{gmFig}[scale=0.42]
  \fill[tmOrange!22] (-4,-4) rectangle (4,4);
  \fill[tmBlue!18] (4,1)--(-1,4)--(-4,-1)--(1,-4)--cycle;
  \draw[Line,line width=0.3pt] (-4,-4) grid (4,4);
  \draw[->,Navy] (-4.6,0)--(4.7,0) node[right,text=Ink]{$x$};
  \draw[->,Navy] (0,-4.6)--(0,4.7) node[above,text=Ink]{$y$};
  \draw[tmOrange!90!black,thick,dashed] (-4,-4) rectangle (4,4);
  \draw[tmBlue,thick] (4,1)--(-1,4)--(-4,-1)--(1,-4)--cycle;
  \foreach \p in {(4,1),(-1,4),(-4,-1),(1,-4)}{\fill[tmBlue] \p circle (3pt);}
  \node[right] at (4,1){$A$}; \node[above] at (-1,4){$B$}; \node[left] at (-4,-1){$C$}; \node[below] at (1,-4){$E$};
  \node[font=\scriptsize] at (2.9,3.2){$5\times3$};
\end{gmFig}
Hình vuông $ABCE$ nằm gọn trong hình vuông lớn có các đỉnh $(\pm4;\,\pm4)$, cạnh $4-(-4)=8$, diện tích $8\cdot8=64$.

Phần thừa gồm $4$ tam giác vuông bằng nhau ở bốn góc. Ví dụ, tam giác ở góc trên bên phải có đỉnh $A(4;\,1)$, $(4;\,4)$, $B(-1;\,4)$: cạnh nằm ngang dài $4-(-1)=5$, cạnh thẳng đứng dài $4-1=3$, diện tích $\dfrac{5\cdot3}{2}=7{,}5$.

Diện tích hình vuông $ABCE$: $64-4\cdot7{,}5=64-30=34$ (đơn vị diện tích).
\begin{tipbox}``Đóng khung rồi trừ bớt'' là cách rất mạnh để tính diện tích một đa giác có các đỉnh nằm ở nút lưới: vẽ hình chữ nhật nhỏ nhất bao quanh, rồi trừ các tam giác vuông (hình thang, hình chữ nhật) ở phía ngoài.\end{tipbox}
\verify{Bốn tam giác vuông bằng nhau (mỗi tam giác $5\times3$) cùng với hình vuông $ABCE$ lấp đầy khung $8\times8$: $4\cdot7{,}5+34=64$. Ngoài ra, cạnh khung $8=5+3$ đúng bằng tổng hai cạnh góc vuông.}
\ans{a) $B(-1;\,4)$, $C(-4;\,-1)$, $E(1;\,-4)$. \ b) Điểm $B$. \ c) $34$ đơn vị diện tích.}
}

\begin{probbox}{35}{\relax}
Điền số thích hợp vào chỗ trống:
\par\smallskip
a) $2{,}4\text{ km}=\ldots\text{ m}$; \qquad b) $75\text{ mm}=\ldots\text{ cm}$; \qquad c) $3\text{ m}^2=\ldots\text{ cm}^2$; \qquad d) $2500\text{ cm}^3=\ldots$ lít.
\par\smallskip
e) Một tấm bảng hình chữ nhật dài $1{,}2$ m, rộng $50$ cm. Tính diện tích tấm bảng theo $\text{cm}^2$ và theo $\text{m}^2$.
\end{probbox}
\probsol{
\textbf{Phân tích đề.} Mỗi câu là một phép đổi đơn vị. Điều cần nhớ: độ dài gấp $10$ lần qua mỗi bậc, diện tích gấp $100$ lần, thể tích gấp $1000$ lần.
\begin{recallbox}$1\text{ km}=1000\text{ m}$; $\;1\text{ cm}=10\text{ mm}$; $\;1\text{ m}^2=10\,000\text{ cm}^2$; $\;1\text{ lít}=1\text{ dm}^3=1000\text{ cm}^3$. Lớn sang nhỏ: nhân; nhỏ sang lớn: chia.\end{recallbox}
\step{1}{Câu a}
km lớn hơn m, nên nhân: $2{,}4\cdot1000=2400$. Vậy $2{,}4\text{ km}=2400\text{ m}$.

\step{2}{Câu b}
mm nhỏ hơn cm, nên chia: $75:10=7{,}5$. Vậy $75\text{ mm}=7{,}5\text{ cm}$.

\step{3}{Câu c}
$3\text{ m}^2=3\cdot10\,000\text{ cm}^2=30\,000\text{ cm}^2$.
\begin{warnbox}Đổi diện tích mà chỉ nhân $100$ (như đổi độ dài m sang cm) sẽ được $300$ -- sai. Hình vuông $1\text{ m}\times1\text{ m}$ là $100\text{ cm}\times100\text{ cm}$, chứa $10\,000$ ô vuông $1\text{ cm}^2$.\end{warnbox}

\step{4}{Câu d}
$2500\text{ cm}^3=2500:1000\text{ dm}^3=2{,}5\text{ dm}^3=2{,}5$ lít.

\step{5}{Câu e -- đưa về cùng đơn vị trước}
Hai kích thước khác đơn vị, nên phải đổi trước khi nhân: $1{,}2\text{ m}=120\text{ cm}$.
Diện tích: $120\cdot50=6000\ (\text{cm}^2)$.
Đổi ra $\text{m}^2$: $6000:10\,000=0{,}6\ (\text{m}^2)$.
\verify{Tính trực tiếp bằng mét: $50\text{ cm}=0{,}5\text{ m}$, nên $1{,}2\cdot0{,}5=0{,}6\ (\text{m}^2)$. Khớp.}
\ans{a) $2400$; \ b) $7{,}5$; \ c) $30\,000$; \ d) $2{,}5$; \ e) $6000\text{ cm}^2=0{,}6\text{ m}^2$.}
}

\begin{probbox}{36}{\stars{1}}
Anna xếp được khối hộp chữ nhật $4\text{ cm}\times3\text{ cm}\times3\text{ cm}$ từ $36$ khối lập phương nhỏ cạnh $1$ cm (như Bài 10). Anna sơn tất cả các mặt bên ngoài của khối hộp, rồi tách rời các khối nhỏ ra. Có bao nhiêu khối nhỏ được sơn đúng $3$ mặt? Đúng $2$ mặt? Đúng $1$ mặt? Không được sơn mặt nào?
\begin{gmFig}[scale=0.55]
  \fill[tmBlue!12] (0,0)--(4,0)--(5.5,1)--(5.5,4)--(1.5,4)--(0,3)--cycle;
  \draw[Ink] (0,0) grid (4,3);
  \foreach \t in {1,2}{\draw[Ink] (4+0.5*\t,0.333*\t)--(4+0.5*\t,3+0.333*\t); \draw[Ink] (0.5*\t,3+0.333*\t)--(4+0.5*\t,3+0.333*\t);}
  \foreach \t in {1,2}{\draw[Ink] (4,\t)--(5.5,\t+1);}
  \foreach \t in {1,2,3}{\draw[Ink] (\t,3)--(\t+1.5,4);}
  \draw[Ink,thick] (0,0) rectangle (4,3); \draw[Ink,thick] (0,3)--(1.5,4)--(5.5,4)--(5.5,1)--(4,0); \draw[Ink,thick] (4,3)--(5.5,4);
\end{gmFig}
\end{probbox}
\probsol{
\textbf{Phân tích đề.} Số mặt được sơn của một khối nhỏ phụ thuộc vào \emph{vị trí} của nó: ở đỉnh, nằm trên cạnh, nằm giữa một mặt, hay nằm hẳn bên trong. Ta đếm từng loại.
\step{1}{Khối ở các đỉnh -- sơn $3$ mặt}
Mỗi khối ở đỉnh của hình hộp lộ ra $3$ mặt (một mặt trước/sau, một mặt trên/dưới, một mặt trái/phải). Hình hộp có $8$ đỉnh, nên có $\mathbf{8}$ khối sơn $3$ mặt.

\step{2}{Khối nằm trên cạnh (trừ hai đầu) -- sơn $2$ mặt}
Hình hộp có $12$ cạnh: $4$ cạnh dài $4$ cm, $4$ cạnh dài $3$ cm (chiều cao) và $4$ cạnh dài $3$ cm (chiều sâu). Trên một cạnh dài $n$ khối, bỏ hai khối ở hai đầu (đã đếm ở bước 1), còn $n-2$ khối.
\begin{itemize}[leftmargin=1.3em,itemsep=0pt,topsep=0pt]
  \item $4$ cạnh dài $4$: mỗi cạnh $4-2=2$ khối, tổng $4\cdot2=8$;
  \item $8$ cạnh dài $3$: mỗi cạnh $3-2=1$ khối, tổng $8\cdot1=8$.
\end{itemize}
Có $8+8=\mathbf{16}$ khối sơn $2$ mặt.

\step{3}{Khối nằm giữa một mặt -- sơn $1$ mặt}
Trên một mặt $m\times n$, bỏ đi viền ngoài, còn $(m-2)\cdot(n-2)$ khối.
\begin{itemize}[leftmargin=1.3em,itemsep=0pt,topsep=0pt]
  \item $2$ mặt $4\times3$ (trước, sau): mỗi mặt $2\cdot1=2$ khối, tổng $4$;
  \item $2$ mặt $4\times3$ (trên, dưới): mỗi mặt $2\cdot1=2$ khối, tổng $4$;
  \item $2$ mặt $3\times3$ (trái, phải): mỗi mặt $1\cdot1=1$ khối, tổng $2$.
\end{itemize}
Có $4+4+2=\mathbf{10}$ khối sơn $1$ mặt.

\step{4}{Khối nằm hẳn bên trong -- không sơn}
Bỏ đi một lớp vỏ ở mỗi phía, phần lõi bên trong có kích thước $(4-2)\times(3-2)\times(3-2)=2\times1\times1$, gồm $\mathbf{2}$ khối.
Tóm tắt kết quả:
\begin{center}\begin{tabular}{|l|c|c|c|c|c|}\hline
Số mặt được sơn & $3$ & $2$ & $1$ & $0$ & Tổng\\\hline
Số khối nhỏ & $8$ & $16$ & $10$ & $2$ & $36$\\\hline
\end{tabular}\end{center}
\begin{parentbox}Nếu có bộ khối lập phương đồ chơi (hoặc khối rubik), hãy cùng con xếp hình hộp $4\times3\times3$, dán giấy màu lên các mặt ngoài rồi tách ra để đếm. Con sẽ tự thấy vì sao khối ở đỉnh có $3$ mặt màu.\end{parentbox}
\verify{Tổng bốn loại phải bằng số khối: $8+16+10+2=36=4\cdot3\cdot3$. Đúng.}
\ans{$3$ mặt: $8$ khối; $2$ mặt: $16$ khối; $1$ mặt: $10$ khối; không sơn: $2$ khối.}
}

\begin{probbox}{37}{\relax}
Một hình chữ nhật có chu vi bằng chu vi của một hình vuông có diện tích $81\text{ cm}^2$. Chiều dài hình chữ nhật gấp đôi chiều rộng. Diện tích hình chữ nhật là
\choices{$81\text{ cm}^2$}{$72\text{ cm}^2$}{$54\text{ cm}^2$}{$36\text{ cm}^2$}{$64\text{ cm}^2$}
\end{probbox}
\probsol{
\textbf{Phân tích đề.} Chuỗi suy luận: diện tích hình vuông $\to$ cạnh hình vuông $\to$ chu vi (chung cho cả hai hình) $\to$ nửa chu vi của hình chữ nhật $\to$ chiều dài, chiều rộng $\to$ diện tích.
\begin{recallbox}Hình vuông cạnh $a$: $S=a\cdot a$, $P=4a$. Hình chữ nhật: $P=2\cdot(\text{dài}+\text{rộng})$, $S=\text{dài}\cdot\text{rộng}$.\end{recallbox}
\step{1}{Cạnh hình vuông}
Cần số $a$ với $a\cdot a=81$: $a=9$ cm (vì $9\cdot9=81$).

\step{2}{Chu vi}
Chu vi hình vuông $=4\cdot9=36$ cm, cũng là chu vi hình chữ nhật.

\step{3}{Tổng chiều dài và chiều rộng}
Nửa chu vi hình chữ nhật: $36:2=18$ cm $=$ dài $+$ rộng.

\step{4}{Gọi ẩn và lập phương trình}
Gọi chiều rộng của hình chữ nhật là $x$ (cm), với $x>0$. Chiều dài gấp đôi chiều rộng nên bằng $2x$ (cm).
Theo bước 3, dài $+$ rộng $=18$, nên ta có phương trình
\[ x+2x=18. \]
Thu gọn: $3x=18$, suy ra $x=18:3=6$.
Vậy chiều rộng là $6$ cm, chiều dài là $2\cdot6=12$ cm.

\step{5}{Diện tích}
$S=12\cdot6=72\ (\text{cm}^2)$.
\begin{tipbox}Hai hình có cùng chu vi $36$ cm, nhưng hình vuông có diện tích $81\text{ cm}^2$ lớn hơn hình chữ nhật ($72\text{ cm}^2$). Trong các hình chữ nhật cùng chu vi, hình vuông có diện tích lớn nhất.\end{tipbox}
\verify{Chu vi hình chữ nhật: $2\cdot(12+6)=36$ cm $=$ chu vi hình vuông; $12=2\cdot6$. Đúng.}
Vậy chọn \textbf{(B)}.
\ans{(B) $72\text{ cm}^2$}
}

\begin{probbox}{38}{\stars{1}}
Một vòng quay may mắn hình tròn được chia thành bốn ô có góc ở tâm lần lượt là $36^\circ$, $x^\circ$, $2x^\circ$ và $3x^\circ$ (hình vẽ).
\begin{enumerate}[label=\alph*),leftmargin=1.6em,itemsep=0pt,topsep=2pt]
  \item Tìm $x$.
  \item Mỗi ô chiếm bao nhiêu phần trăm diện tích vòng quay?
  \item Khi quay, kim \emph{không} dừng ở ô lớn nhất trong khoảng bao nhiêu phần trăm số lần?
\end{enumerate}
\begin{gmFig}[scale=0.9]
  \fill[tmPurple!20] (0,0)--(90:1.7) arc[start angle=90,end angle=126,radius=1.7]--cycle;
  \fill[tmBlue!20] (0,0)--(126:1.7) arc[start angle=126,end angle=180,radius=1.7]--cycle;
  \fill[tmGreen!20] (0,0)--(180:1.7) arc[start angle=180,end angle=288,radius=1.7]--cycle;
  \fill[tmOrange!25] (0,0)--(288:1.7) arc[start angle=288,end angle=450,radius=1.7]--cycle;
  \draw[Ink,thick] (0,0) circle (1.7);
  \foreach \a in {90,126,180,288}{\draw[Ink,thick] (0,0)--(\a:1.7);}
  \node[font=\footnotesize] at (108:1.2){$36^\circ$};
  \node[font=\footnotesize] at (153:1.1){$x^\circ$};
  \node[font=\footnotesize] at (234:1.0){$2x^\circ$};
  \node[font=\footnotesize] at (9:0.95){$3x^\circ$};
\end{gmFig}
\end{probbox}
\probsol{
\textbf{Phân tích đề.} Bốn góc ở tâm ghép lại thành một vòng trọn. Từ đó lập được một đẳng thức để tìm $x$, rồi đổi từng góc ra phần trăm.
\begin{recallbox}Các góc ở tâm quanh một điểm có tổng $360^\circ$. Ô có góc ở tâm $n^\circ$ chiếm $\dfrac{n}{360}$ hình tròn.\end{recallbox}
\step{1}{Câu a -- lập đẳng thức}
$36+x+2x+3x=360$, tức $36+6x=360$.
Bớt $36$ ở cả hai vế: $6x=324$. Chia cho $6$: $x=54$.

\step{2}{Câu b -- các góc và phần trăm}
Bốn góc: $36^\circ$, $54^\circ$, $108^\circ$, $162^\circ$.
\[ \frac{36}{360}=10\%,\qquad \frac{54}{360}=15\%,\qquad \frac{108}{360}=30\%,\qquad \frac{162}{360}=45\%. \]
(Mẹo tính: $\frac{36}{360}=\frac1{10}$; các góc còn lại gấp $1{,}5$; $3$; $4{,}5$ lần $36^\circ$, nên phần trăm lần lượt là $15\%$, $30\%$, $45\%$.)

\step{3}{Câu c -- dùng phần bù}
Ô lớn nhất là ô $162^\circ$, chiếm $45\%$. Kim \emph{không} dừng ở ô đó trong khoảng $100\%-45\%=55\%$ số lần.
\begin{warnbox}Kết quả ``$55\%$'' chỉ là \emph{khoảng} số lần sau rất nhiều lượt quay. Trong vài lượt quay, số lần thực tế có thể khác; tỉ lệ diện tích cho biết kim dừng ở đâu \emph{thường xuyên đến mức nào} về lâu dài.\end{warnbox}
\verify{$36+54+108+162=360$ và $10\%+15\%+30\%+45\%=100\%$.}
\ans{a) $x=54$. \ b) $10\%$, $15\%$, $30\%$, $45\%$. \ c) Khoảng $55\%$.}
}

\begin{probbox}{39}{\stars{1}}
Một hình chữ nhật dài $15$ cm, rộng $6$ cm. Trên cạnh trên, lấy hai đoạn ở hai đầu có cùng độ dài $a$ cm rồi cắt bỏ hai tam giác vuông ở hai góc trên, còn lại hình thang tô màu như hình. Biết diện tích hình thang là $72\text{ cm}^2$. Tìm $a$ và độ dài đáy nhỏ của hình thang.
\begin{gmFig}[scale=0.42]
  \fill[tmBlue!35] (0,0)--(15,0)--(12,6)--(3,6)--cycle;
  \draw[Ink,thick] (0,0) rectangle (15,6);
  \draw[Ink,thick] (0,0)--(3,6); \draw[Ink,thick] (15,0)--(12,6);
  \node[above] at (1.5,6){$a$}; \node[above] at (13.5,6){$a$};
  \node[below] at (7.5,0){$15$ cm}; \node[right] at (15,3){$6$ cm};
\end{gmFig}
\end{probbox}
\probsol{
\textbf{Phân tích đề.} Bài 14 cho kích thước và hỏi diện tích; bài này làm ngược lại: cho diện tích, hỏi kích thước. Ta viết diện tích hình thang theo $a$ bằng cách ``hình chữ nhật trừ hai tam giác'', rồi tìm $a$.
\begin{recallbox}$S_{\text{hình chữ nhật}}=$ dài $\times$ rộng; $\;S_{\text{tam giác vuông}}=\dfrac{\text{tích hai cạnh góc vuông}}{2}$.\end{recallbox}
\step{1}{Diện tích hình chữ nhật}
$15\cdot6=90\ (\text{cm}^2)$.

\step{2}{Tổng diện tích hai tam giác bị cắt}
Diện tích hai tam giác $=90-72=18\ (\text{cm}^2)$, nên mỗi tam giác có diện tích $18:2=9\ (\text{cm}^2)$.

\step{3}{Tìm $a$}
Mỗi tam giác vuông có hai cạnh góc vuông là $a$ và $6$:
\[ \frac{a\cdot6}{2}=9\quad\Rightarrow\quad 3\cdot a=9\quad\Rightarrow\quad a=3. \]

\step{4}{Đáy nhỏ của hình thang}
Đáy nhỏ $=15-a-a=15-3-3=9$ cm.
\verify{Theo công thức hình thang: $\dfrac{(15+9)\cdot6}{2}=\dfrac{24\cdot6}{2}=72\ (\text{cm}^2)$. Đúng.}
\begin{tipbox}Cũng có thể giải bằng công thức hình thang: đáy nhỏ $b$ thoả mãn $\dfrac{(15+b)\cdot 6}{2}=72$, tức $(15+b)\cdot3=72$, $15+b=24$, $b=9$; rồi $a=(15-9):2=3$.\end{tipbox}
\ans{$a=3$ cm; đáy nhỏ $9$ cm.}
}

\begin{probbox}{40}{\stars{1}}
Tam giác $ABC$ vuông tại $A$ có $AB=6$ cm, $AC=8$ cm và $BC=10$ cm.
\begin{enumerate}[label=\alph*),leftmargin=1.6em,itemsep=0pt,topsep=2pt]
  \item Tính diện tích tam giác $ABC$.
  \item Tính độ dài đường cao $AH$ kẻ từ $A$ xuống cạnh $BC$.
  \item Gọi $M$ là trung điểm của $BC$. Tính diện tích tam giác $ABM$.
\end{enumerate}
\begin{gmFig}[scale=0.55]
  \coordinate (B) at (0,0); \coordinate (C) at (10,0); \coordinate (A) at (3.6,4.8);
  \coordinate (H) at (3.6,0); \coordinate (M) at (5,0);
  \fill[tmBlue!10] (A)--(B)--(C)--cycle;
  \draw[Ink,thick] (A)--(B)--(C)--cycle;
  \draw[AnsRed,thick,dashed] (A)--(H); \draw[Ink] (H) rectangle ++(0.3,0.3);
  \draw[tmGreen!70!black,thick] (A)--(M);
  \draw[Ink] ($(A)+(-0.24,-0.32)$)--($(A)+(0.08,-0.56)$)--($(A)+(0.32,-0.24)$);
  \node[above] at (A){$A$}; \node[below left] at (B){$B$}; \node[below right] at (C){$C$};
  \node[below] at (H){$H$}; \node[below] at (M){$M$};
  \node[above left=1pt] at (1.8,2.4){$6$}; \node[above right=1pt] at (6.8,2.4){$8$};
\end{gmFig}
\end{probbox}
\probsol{
\textbf{Phân tích đề.} Một tam giác có thể chọn \emph{bất kì} cạnh nào làm đáy, miễn là dùng đúng chiều cao ứng với đáy đó. Diện tích luôn như nhau. Ta tính diện tích theo cách dễ (đáy $AB$, chiều cao $AC$), rồi dùng nó để tìm chiều cao ứng với đáy $BC$.
\begin{recallbox}$S=\dfrac{\text{đáy}\cdot\text{chiều cao}}{2}$, suy ra chiều cao $=\dfrac{2\cdot S}{\text{đáy}}$. Tam giác vuông: hai cạnh góc vuông là đáy và chiều cao của nhau.\end{recallbox}
\step{1}{Câu a -- chọn đáy $AB$}
Vì góc $A$ vuông nên $AC$ vuông góc với $AB$: lấy $AB$ làm đáy thì $AC$ là chiều cao.
\[ S_{ABC}=\frac{6\cdot8}{2}=24\ (\text{cm}^2). \]

\step{2}{Câu b -- chọn đáy $BC$}
Lấy $BC=10$ cm làm đáy, chiều cao tương ứng là $AH$. Diện tích vẫn là $24\text{ cm}^2$:
\[ \frac{10\cdot AH}{2}=24\quad\Rightarrow\quad 10\cdot AH=48\quad\Rightarrow\quad AH=4{,}8\text{ cm}. \]

\step{3}{Câu c -- trung điểm chia đôi diện tích}
\begin{recallbox}Hai tam giác có đáy bằng nhau và chung chiều cao thì có diện tích bằng nhau.\end{recallbox}
Hai tam giác $ABM$ và $ACM$ có đáy $BM=MC=5$ cm (vì $M$ là trung điểm) và chung chiều cao $AH$ (cùng kẻ từ $A$ xuống đường thẳng $BC$). Vậy chúng có diện tích bằng nhau, mỗi tam giác bằng nửa tam giác $ABC$:
\[ S_{ABM}=\frac{5\cdot4{,}8}{2}=12\ (\text{cm}^2)=\frac{24}{2}. \]
\begin{warnbox}Đừng lấy ``đáy $BC$ nhân cạnh $AB$'' -- cạnh $AB$ không vuông góc với $BC$, nên không phải chiều cao ứng với đáy $BC$.\end{warnbox}
\verify{$AH=4{,}8$ nhỏ hơn cả $AB=6$ và $AC=8$ -- hợp lí, vì đường vuông góc là đường ngắn nhất từ $A$ đến đường thẳng $BC$. Và $\frac{10\cdot4{,}8}{2}=24$.}
\ans{a) $24\text{ cm}^2$. \ b) $AH=4{,}8$ cm. \ c) $12\text{ cm}^2$.}
}
